\scanpage{32}
% archive sheet number: 31
\diagram{\begin{tikzcd}[ampersand replacement=\&]
e \arrow[r] \& G_m^{\text{tordu}} \arrow[r] \& G \arrow[r] \& Z/2 \arrow[r] \& e\\
e \arrow[r] \& \mathbb Z/3 \arrow[u]\arrow[r] \& \mathfrak{S}_3 \arrow[r] \& \mathbb Z/2 \arrow[u,equal]\arrow[r] \& e
\end{tikzcd}}

On suppose tjrs qu'il n'y a pas de \unsure{pt|pts} $s$ de car. 2 ou 3 dans $S$. Alors \deleted{\unsure{$\mathfrak S_3$|$G$|le} \unsure{opère|groupe}} $\mathfrak S_3\times Z_2$ opère ``tamely'' sur $P^1_S$ [en effet, il est compatible avec la fibration $P^1_S\to S$, et opère\unsure{'|} ``tamely'' sur les fibres $\ldots$] Le quotient \inserted{par \unsure{ce|un|le} groupe ou un ss-groupe} \deleted{s'identifie} \inserted{est isomorphe} encore à $P^1_S$ [car il est \uline{\unsure{lisse|simple}} sur $S$, et rationnel sur $S$ d'après Lüroth]. Les schémas de pts fixes $(0,1,\infty)$ $(\unsure{2|3},-1,1/2)$ $(\bar S=x',x'')$ -- qui sont disjoints, et chacun stables sous $\mathfrak S_3$, qui \deleted{y} \inserted{sur chacun} opère de façon \deleted{\unsure{simplement}} transitive sur chaque fibre; donc on obtient dans $P^1_S/\mathfrak S_3$ trois \uline{sections} distinctes [à expliciter $\ldots$] qu'on peut utiliser pour identifier $P^1_S/\mathfrak S_{\unsure{3|2}}$ \unsure{.|à} $P^1_S$ avec les sections $\infty,0,1$. De $=$, \deleted{de se donner des \unsure{sect.|sect}} $P^1_S/(\mathbb Z/3)$ est muni d'abord de deux sections qui correspondent à $(0,1,\infty)$ et $(2,-1,1/2)$, et qui

